% 电路与电子学III 课程笔记 —— 第十二章（教材第12章）非线性电阻电路
\chapter{非线性电阻电路}

\section{非线性电阻元件的特性}

含有非线性元件的电路，其电路方程是非线性的，这种电路称为\textbf{非线性电路}。非线性电路广泛存在：
一方面是某些元件不可避免的非线性特性（如二极管、晶体管），另一方面是人为设计的非线性电路，
用以实现线性电路无法实现的功能。分析非线性电路要比分析线性电路复杂得多，
\textbf{基于线性方程的各种电路定理不能直接用于非线性电路}。作为基础，本章研究最简单的非线性电路
——非线性电阻电路。

\subsection{定义与典型特性}
非线性电阻元件的\textbf{约束方程是非线性的}：其电压、电流关系不是通过 $u-i$ 平面坐标原点的直线，
因而不满足欧姆定律，只能写成
\begin{equation}
  u=f(i)\qquad\text{或}\qquad i=g(u)
\end{equation}
例如半导体 PN 结二极管，根据半导体物理可得其电压、电流的解析关系为
\begin{keyformula}{i=I_S\left(\mathrm{e}^{u/U_T}-1\right)}
\end{keyformula}
式中 $I_S$ 为反向饱和电流，$U_T$ 为温度电压当量，二者在二极管型号与工作温度一定时都是常量
（室温下 $U_T\approx26\ \mathrm{mV}$）。加正向电压时特性基本按指数规律上升；加反向电压时
$\mathrm{e}^{u/U_T}$ 很快趋于零，电流近似等于 $-I_S$（很小的量）。任何元件的数学模型都有适用范围：
当反向电压超过反向击穿电压（$|u|>|U_{BR}|$）时，反向电流急剧增加，不再遵循该解析式。

实际上多数非线性电阻元件难以从理论上得出解析表达式，通常只能用\textbf{曲线或数据表}描述。

\subsection{按控制量分类}
\begin{itemize}
  \item \textbf{单调型}：$u$ 与 $i$ 互为单值函数（如 PN 结二极管），$u=u(i)$ 与 $i=i(u)$ 都可单值表示。
  \item \textbf{流控型}（current controlled）：$u$ 是 $i$ 的单值函数，记作 $u=u(i)$；而同一电压下电流可能出现
        多值，即其反函数是多值函数。典型例子是辉光管，其特性在第 I 象限呈“S”形。
  \item \textbf{压控型}（voltage controlled）：$i$ 是 $u$ 的单值函数，记作 $i=i(u)$；而同一电流下电压可能多值。
        典型例子是隧道二极管，其特性在第 I 象限呈“N”形。
\end{itemize}
区分单调型、压控型和流控型的主要原因是：\textbf{对不同类型要选用不同的方法列写电路方程}
（流控型取 $i$ 为变量，压控型取 $u$ 为变量）。

\subsection{双向性与方向性}
线性电阻没有方向性，其特性曲线对称于坐标原点，具有\textbf{双向性}，使用时不必考虑极性。
有些非线性电阻（如辉光管）的特性仍是双向的，即 $-u=u(-i)$；而 PN 结二极管、隧道二极管
\textbf{有方向性}，正向与反向的导电性不同，特性曲线对原点不对称，电路符号上有明显的方向标志，
使用时应按方向正确连接。

\subsection{静态电阻与动态电阻}
非线性电阻的电阻值不是常数，随工作点而变，必须区分两种“电阻”：
\begin{keyformula}{R=\frac{u}{i}\ \ (\text{\songti 静态电阻}),\qquad
  R_d=\frac{\mathrm{d}u}{\mathrm{d}i}\ \ (\text{\songti 动态电阻})}
\end{keyformula}
静态电阻 $R$ 是工作点与原点连线的斜率的倒数，反映该工作点上的整体伏安关系；
动态电阻 $R_d$（又称增量电阻、微变电阻）是特性曲线在该点切线的斜率，反映工作点附近的局部变化关系。
同理有静态电导 $G=i/u$ 与动态电导 $G_d=\mathrm{d}i/\mathrm{d}u=1/R_d$。
$R$ 恒为正，而 $R_d$ \textbf{可以为负}（S 形、N 形特性的下降段），负动态电阻向电路提供能量。

\begin{corepoint}{静态电阻与动态电阻一般\textbf{不相等}，只有线性电阻（特性为过原点的直线）二者才相等。
易错点：小信号分析中用的是\textbf{动态电阻} $R_d$，而算工作点功率、效率时用\textbf{静态电阻}（即 $U_Q/I_Q$）。
图中“切线斜率”的取法：$R_d=\mathrm{d}u/\mathrm{d}i$，斜率越大动态电阻越大，切勿取倒数。}
\end{corepoint}

\section{非线性直流电路方程}

电路方程决定于\textbf{结构约束}和\textbf{支路约束}：前者即基尔霍夫定律方程，它们不涉及元件性质，
总是线性代数方程；后者取决于支路中的元件，若存在非线性电阻，就会出现非线性代数方程。因此
\textbf{叠加定理、齐性原理、戴维南定理、诺顿定理等不能直接用于解非线性电路}。

实际电路中常是线性元件与非线性元件同时存在，往往可以先把线性部分尽量划分出来加以等效化简，
以便建立尽可能简单的非线性电路方程。

\subsection{只含一个非线性电阻时的求解步骤}
若电路只含一个非线性电阻且主要关心它的解答，可把电路划分为\textbf{线性部分}和\textbf{非线性部分}：
\begin{enumerate}
  \item 利用线性电路的戴维南定理（或诺顿定理）对线性部分等效化简，得到只含一个电压源 $U_{OC}$、
        一个等效电阻 $R$ 与非线性电阻串联的简单电路；
  \item 列写该电路方程：若为\textbf{流控型}（$u=u(i)$），以电流 $i$ 为变量列 KVL
        \begin{equation}
          Ri+u(i)=U_{OC}
          \label{eq:ch12-flow}
        \end{equation}
        若为\textbf{压控型}（$i=i(u)$），以电压 $u$ 为变量列 KVL
        \begin{equation}
          Ri(u)+u=U_{OC}
          \label{eq:ch12-volt}
        \end{equation}
  \item 求解式 \eqref{eq:ch12-flow} 或 \eqref{eq:ch12-volt}，得到非线性电阻上的解答 $(U,I)$；
  \item 若还要进一步求线性部分的解答，可用电压源或电流源\textbf{置换}非线性电阻（置换定理），
        得到线性直流电路再求解。
\end{enumerate}

\subsection{含多个非线性电阻时的方程列写}
若电路含有较多的非线性电阻，或需要计算较多的电压、电流，则宜对电路列写联立方程组：
用节点电压法或回路电流法把结构约束写成线性方程，再把各非线性电阻的支路约束代入，
得到\textbf{非线性代数方程组}。这类方程组一般没有解析解，需用下一节的数值方法求解。

\begin{corepoint}{非线性电路的三个关键差别：\textbf{（1）}不能叠加，故“分别求各电源作用再相加”失效；
\textbf{（2）}解可能不唯一（多个工作点），必须逐个检验其物理合理性（如功率、是否落在特性曲线的适用段内）；
\textbf{（3）}“线性部分可用戴维南定理等效”与“戴维南定理可用于非线性电路”是两回事——前者成立，
因为等效只涉及线性部分；后者不成立。}
\end{corepoint}

\section{数值分析法}

求解非线性电路归结为求非线性方程（组）的根，一般不易求解析解，常用数值分析法，
其中最基本的是\textbf{牛顿—拉夫逊法}（切线法）。

\subsection{牛顿—拉夫逊迭代公式}
设要求解方程 $f(x)=0$，取初值 $x_0$，迭代公式为
\begin{keyformula}{x_{k+1}=x_k-\frac{f(x_k)}{f'(x_k)},\qquad k=0,1,2,\cdots}
\end{keyformula}
其几何意义是：在点 $(x_k,f(x_k))$ 处作曲线 $f(x)$ 的切线，以切线与横轴的交点作为下一次的近似值 $x_{k+1}$。
当 $|x_{k+1}-x_k|$ 小于规定容许误差 $\varepsilon$ 时停止迭代，取 $x_{k+1}$ 为解。

\subsection{应用于非线性电阻电路}
以含二极管的电路为例，回路电压方程为
\begin{equation}
  u_R+u-U_S=RI+u-U_S=0
\end{equation}
将二极管的特性 $i=I_S(\mathrm{e}^{u/U_T}-1)$ 代入，得到关于 $u$ 的非线性方程
\begin{equation}
  f(u)=RI_S\left(\mathrm{e}^{u/U_T}-1\right)+u-U_S=0
\end{equation}
其导数 $f'(u)=R\dfrac{I_S}{U_T}\mathrm{e}^{u/U_T}+1$，代入迭代公式即可逐次逼近。若特性以多项式给出，
也可直接求多项式的根（如 MATLAB 的 \texttt{roots}）。

\subsection{收敛性与注意点}
牛顿法收敛快（平方收敛），但有条件：\textbf{初值须在真解附近}，且迭代过程中 $f'(x_k)\ne0$；
若初值选得不好或曲线有拐点，迭代可能发散或振荡。改善办法：选取合理初值（如用图解法定出大致位置）、
限制每次修正量（阻尼）、改用弦截法或二分法等。

\begin{corepoint}{数值法的答案必须\textbf{回代检验}：把求得的 $(U,I)$ 代回原方程，残差应小于容许误差。
另外多解时不能只报第一个收敛值——非线性电路常有两个甚至三个工作点，需按物理意义（如
$u$ 是否在特性曲线适用范围内、功率是否为负）取舍。}
\end{corepoint}

\section{分段线性近似法}

非线性电阻元件的特性通常可以测定其电压、电流特性曲线，但有时不能找出它的解析表达式。
此时可用若干段直线去逼近该曲线，每一段直线对应一个线性方程，从而把非线性问题化为一系列线性问题，
这种方法称为\textbf{分段线性近似法}。

\subsection{折线的数学表示}
把特性曲线分成若干区段，第 $k$ 段用直线近似：
\begin{keyformula}{u=R_{dk}i+U_{0k}\qquad\text{或}\qquad i=G_{dk}u+I_{0k}}
\end{keyformula}
其中 $R_{dk}=\Delta u/\Delta i$ 是该段的\textbf{动态电阻}（斜率的倒数），$U_{0k}$ 是该段直线在电压轴上的
截距。因此每一段都可等效为一个动态电阻与一个电压源（或动态电导与一个电流源）的串联（并联）组合。

\subsection{求解步骤}
\begin{enumerate}
  \item 把每个非线性电阻的特性曲线用折线近似，确定各段的 $R_{dk}$、$U_{0k}$；
  \item 对每个非线性电阻\textbf{假定它工作在某一区段}，用该段的线性模型替代它，得到相应的线性电路；
  \item 解这个线性电路，求出各非线性电阻的电压、电流；
  \item \textbf{检验}：所得工作点是否落在所假定的区段内。若全部符合，则该解有效；
        否则换一种区段组合重新计算。
\end{enumerate}
当电路含多个非线性电阻时，区段组合有多种，只有各非线性电阻的解都落在各自假定的区段内，
该组合才是自洽的；因此可能得到多组自洽解，即多个工作点。

\subsection{特点}
优点是每一段都是线性电路，可直接搬用线性电路的一切方法（等效、节点法、回路法），且能处理多值特性、
便于求多个工作点；缺点是折线近似带来误差，区段组合多时试算量大。

\begin{corepoint}{分段线性法的核心是\textbf{“假定—求解—校验”}三步闭环，校验一步不能省：
算出的工作点若越出假定区段，说明该解无效，必须换组合重算。常见错误是算出一个答案就直接写结论，
而该点其实落在折线的另一段上。}
\end{corepoint}

\section{图解法}

当非线性电阻的特性以曲线给出时，可直接用作图法求工作点，物理概念直观。

\subsection{曲线相交法（负载线法）}
把非线性电阻以外的线性部分等效为 $U_{OC}$ 与 $R$ 串联，其外特性为
\begin{equation}
  u=U_{OC}-Ri
\end{equation}
这是一条斜率为 $-R$、在电压轴上截距为 $U_{OC}$ 的直线，称为\textbf{负载线}。把它与非线性电阻的
$u-i$ 特性曲线画在同一坐标系中，\textbf{两曲线的交点即为工作点} $(U,I)$：交点同时满足线性部分的外特性
与非线性元件的特性。

\subsection{曲线相加法}
当电路含多个非线性电阻时，可用曲线相加求等效特性：
\begin{itemize}
  \item \textbf{串联}：各元件电流相同，故在同一电流下把各元件的电压相加，得等效特性 $u(i)=u_1(i)+u_2(i)$；
  \item \textbf{并联}：各元件电压相同，故在同一电压下把各元件的电流相加，得等效特性 $i(u)=i_1(u)+i_2(u)$。
\end{itemize}
求出等效特性后，再与线性部分的负载线相交，即得工作点；必要时逐级回代求各元件的电压、电流。

\subsection{适用范围}
图解法适用于特性曲线已知且（在所讨论区段内）为单值的情形，对单调型、流控型、压控型都可使用，
只需注意相加时选对公共变量。其缺点是精度受作图与读数限制，多解时需仔细辨认各交点。

\begin{corepoint}{作图法的两个易错点：\textbf{（1）}串联相加时保持\textbf{电流}相同、把\textbf{电压}相加；
并联相加时保持\textbf{电压}相同、把\textbf{电流}相加——公共变量选错则曲线全错。
\textbf{（2）}负载线的斜率是 $-R$ 而不是 $R$，截距是 $U_{OC}$（开路电压），不要与短路电流混淆。}
\end{corepoint}

\section{小信号分析法}

工程上常见这样的情况：非线性电阻上除直流偏置（大信号）外，还叠加了一个幅度很小的交流信号。
此时可在直流工作点附近把特性\textbf{线性化}，用小信号分析法求解。

\subsection{基本原理}
设直流电源单独作用时，非线性电阻工作在\textbf{静态工作点} $Q(U_Q,I_Q)$；再叠加小信号 $\Delta u_S$ 后，
工作点在其附近微小变动。把特性 $i=i(u)$ 在 $Q$ 点作泰勒展开并只取一阶项：
\begin{equation}
  i\approx i(U_Q)+\left.\frac{\mathrm{d}i}{\mathrm{d}u}\right|_{Q}\Delta u
  =I_Q+G_d\Delta u
\end{equation}
即小信号分量之间满足线性关系
\begin{keyformula}{\Delta i=G_d\Delta u,\qquad \Delta u=R_d\Delta i,\qquad
  R_d=\left.\frac{\mathrm{d}u}{\mathrm{d}i}\right|_{Q}}
\end{keyformula}
其中 $R_d$ 是工作点 $Q$ 处的\textbf{动态电阻}。

\subsection{小信号等效电路与解题步骤}
\begin{enumerate}
  \item 令小信号源为零，求出直流工作点 $Q(U_Q,I_Q)$（用前几节的方法）；
  \item 在 $Q$ 点求动态电阻 $R_d=\left.\dfrac{\mathrm{d}u}{\mathrm{d}i}\right|_{Q}$（由特性曲线切线或解析式求导）；
  \item 画\textbf{小信号等效电路}：原电路拓扑不变，直流电源置零（电压源短路、电流源开路），
        非线性电阻用 $R_d$ 替代，只保留小信号源；
  \item 用线性电路方法求小信号响应 $\Delta u$、$\Delta i$；
  \item 叠加得全解：$u=U_Q+\Delta u$，$i=I_Q+\Delta i$。
\end{enumerate}

\subsection{适用条件}
小信号分析法成立的条件是：\textbf{（1）}特性在 $Q$ 点附近连续可导；\textbf{（2）}小信号幅度足够小，
使二阶及以上的高阶项可以忽略。它把非线性问题化为“直流工作点 $+$ 线性小信号”两步，是电子线路
（放大器、二极管检波等）分析的基本方法。

\begin{corepoint}{小信号等效电路中\textbf{直流电源必须置零}、非线性电阻必须换成\textbf{动态电阻}
（而不是静态电阻 $U_Q/I_Q$），这是两个最常见的错误。若特性下降段 $R_d<0$，等效电路中就出现负电阻，
这并非算错。最后回答问题时要把静态量与信号量分别写清，不可混为一个表达式。}
\end{corepoint}

\section{本章小结}
\begin{enumerate}
  \item 非线性电阻的电压、电流关系不满足欧姆定律，其约束方程是非线性的（解析式、曲线或数据表）。
        按控制量分为单调型、流控型（$u=u(i)$ 单值，如辉光管）和压控型（$i=i(u)$ 单值，如隧道二极管）；
        按对称性分为双向性与有方向性的元件。
  \item 非线性电阻有静态电阻 $R=u/i$ 与动态电阻 $R_d=\mathrm{d}u/\mathrm{d}i$ 之分，后者可为负值。
  \item 非线性电路的结构约束仍是线性方程，但支路约束是非线性的，故整体方程为非线性代数方程；
        叠加定理、齐性原理、戴维南定理等不能直接使用，但可对\textbf{线性部分}作等效化简。
  \item 只含一个非线性电阻时，可按“划分线性/非线性—线性部分等效—按控制类型列方程—求解—置换求其余量”
        的步骤处理。
  \item 数值分析法以牛顿—拉夫逊法为代表：$x_{k+1}=x_k-f(x_k)/f'(x_k)$，收敛快但对初值敏感，
        结果须回代检验。
  \item 分段线性近似法用折线代替曲线，每段 $u=R_{dk}i+U_{0k}$，按“假定区段—解线性电路—校验”闭环进行，
        可求出多个工作点。
  \item 图解法用负载线 $u=U_{OC}-Ri$ 与特性曲线求交点；多非线性电阻时用曲线相加（串联加电压、
        并联加电流）求等效特性。
  \item 小信号分析法在工作点处线性化：$\Delta u=R_d\Delta i$，$R_d$ 为 $Q$ 点的动态电阻；
        小信号等效电路中直流电源置零、非线性电阻用 $R_d$ 替代，全解为静态量与信号量之和。
\end{enumerate}

\endinput
